% Build: pdflatex Architecture.tex, twice (table of contents).
\documentclass[11pt,a4paper]{article}
\usepackage[T1]{fontenc}
\usepackage[utf8]{inputenc}
\usepackage{lmodern}
\usepackage{microtype}
\usepackage{parskip}
\setlength{\emergencystretch}{2em}
\usepackage{amsmath,amssymb,mathtools}
\usepackage{booktabs,array}
\usepackage[a4paper,margin=25mm,headheight=14pt]{geometry}
\usepackage{fancyhdr}
\usepackage{xcolor}
\usepackage[colorlinks=true,linkcolor=blue!55!black,urlcolor=blue!55!black,citecolor=blue!55!black]{hyperref}

\pagestyle{fancy}
\fancyhf{}
\lhead{CASC--NTP \textbullet{} Configuration A}
\rhead{Mathematical architecture}
\cfoot{\thepage}
\renewcommand{\headrulewidth}{0.4pt}

\newcommand{\R}{\mathbb{R}}
\newcommand{\rmsn}{\operatorname{RMSNorm}}
\newcommand{\SiLU}{\operatorname{SiLU}}
\newcommand{\softplus}{\operatorname{softplus}}
\newcommand{\softmax}{\operatorname{softmax}}
\newcommand{\gbar}{\bar g}
\newcommand{\Gbar}{\bar G}
\newcommand{\up}{{\uparrow}}
\newcommand{\dn}{{\downarrow}}
\newcommand{\one}{\mathbf{1}}
\newcommand{\op}{\operatorname{op}}
\newcommand{\T}{\top}

\title{\textbf{CASC--NTP}\\[6pt]
\large Mathematical Architecture of the Geometric Cascade Machine\\[2pt]
\normalsize Configuration A: $F$--$B$--$M$--$M$--$V$--$M$--$S$}
\author{CASC Project}
\date{Architecture specification \textbullet{} October 2026}

\begin{document}
\maketitle
\thispagestyle{plain}

\begin{abstract}
CASC--NTP is an autoregressive neural architecture with a recurrent causal carrier, a communication
attention mechanism, a separate control attention mechanism, and a sequence of learned low-rank
geometric operators. Each operator stage acts in a lifted operational space but returns its update to
the hidden state before the next stage is evaluated. The next stage therefore conditions on the
representation already transformed by previous stages. This document specifies the \emph{machine
itself}: its parameterization, tensor dimensions, recurrence, routing simplices, operator families,
composition law, and lexical readout. A final section records the optimization mechanism used in the
documented implementation: standard reverse-mode automatic differentiation (backpropagation) of
next-token cross-entropy, followed by AdamW.
\end{abstract}

\tableofcontents
\newpage

% =====================================================================================================
\section{Scope and mathematical conventions}

The specified machine is \textbf{CASC--NTP, Configuration A}. It has four structurally identical blocks
with independent trainable parameters. Each block processes a causal sequence, communicates recurrent
memory into the hidden state, then applies seven consecutive geometric updates of families
\begin{equation}\label{eq:families}
(f_1,\dots,f_7)=(F,B,M,M,V,M,S).
\end{equation}
The order is architectural: the three occurrences of $M$ have their own operator parameters and cannot
be treated as one shared metric operator.

Throughout, activations are \emph{row vectors}; learned matrices multiply on the right. The position
index is $i\in\{1,\dots,L\}$, block index is $b\in\{1,\dots,S\}$, geometric stage is
$k\in\{1,\dots,K\}$, head index is $a\in\{1,\dots,H\}$, and operator slot is $j\in\{0,\dots,5\}$.
Unless otherwise specified, all formulas for a block omit the superscript $(b)$ to keep the notation
readable. Operators and parameters in distinct blocks are independent.

\subsection{Dimensions}

\begin{center}
\begin{tabular}{p{0.15\textwidth}p{0.47\textwidth}r}
\toprule
Symbol & Meaning & Configuration A\\
\midrule
$V$ & Vocabulary size & 64{,}000\\
$d_x$ & Hidden and lexical embedding width & 1{,}024\\
$d_g$ & Recurrent carrier / attention width & 512\\
$H$ & Number of heads in each attention & 8\\
$d_h$ & Width of each attention head ($Hd_h=d_g$) & 64\\
$D$ & Lifted operational dimension & 4{,}096\\
$r$ & Local operator rank & 32\\
$K$ & Geometric stages per block & 7\\
$S$ & Independent blocks & 4\\
$L$ & Causal sequence length & Configurable\\
\bottomrule
\end{tabular}
\end{center}

Sequence length $L$ is \emph{not} fixed by the architecture. The training code defaults to $L=512$. The published run used $L=96$ for its initial training phase and an extension to $L=384$ near the end. These are protocol choices, not different definitions of the
machine.

\subsection{Elementary functions and normalizations}

For scalar or elementwise arguments, define
\begin{equation}
\sigma(t)=\frac{1}{1+e^{-t}},\qquad \softplus(t)=\log(1+e^{t}),\qquad \SiLU(t)=t\sigma(t).
\end{equation}
Hadamard multiplication is written $\odot$. The numerical stabilization constant is $\varepsilon=10^{-6}$.
For $x\in\R^d$, the parameter-free root-mean-square normalization is
\begin{equation}
\operatorname{rms}(x)=\frac{x}{\sqrt{\|x\|_2^2/d+\varepsilon}}.
\end{equation}
The trainable version uses a learned coordinatewise gain $\gamma\in\R^d$:
\begin{equation}
\rmsn_d(x)=\gamma\odot\operatorname{rms}(x).
\end{equation}
Different uses of RMSNorm have their own scale parameters unless parameter sharing is explicitly
specified. We use
\begin{equation}
\Delta^{n-1}=\Bigl\{\alpha\in\R^n_{\ge 0}:\ \sum_{j=1}^{n}\alpha_j=1\Bigr\}
\end{equation}
for the $(n-1)$-dimensional probability simplex.

% =====================================================================================================
\section{Lexical interface and block composition}

Let $y_1,\dots,y_L\in\{0,\dots,V-1\}$ be token indices. The trainable embedding is
\begin{equation}
E\in\R^{V\times d_x},\qquad E_{q\ell}\sim\mathcal N(0,0.02^2),\qquad x_i^{[0]}=E[y_i]\in\R^{d_x}.
\end{equation}
The same matrix is used by the final language-model head (weight tying).

With $B_b$ denoting the complete transformation performed by block $b$, the machine body is
\begin{equation}
x^{[S]}=B_S\circ B_{S-1}\circ\cdots\circ B_1\bigl(x^{[0]}\bigr),\qquad S=4.
\end{equation}
Every block is internally divided into a \emph{recurrent memory construction}, a \emph{causal
communication read}, and a \emph{seven-stage controlled geometric cascade}. Shared parameters are
shared \emph{across the seven stages within a block}, not between distinct blocks.

% =====================================================================================================
\section{Causal memory: signed-simplex GRU}

\subsection{Memory input and recurrent gates}

At position $i$ of a fixed block, write $x_i\in\R^{d_x}$ for the block input. Project it into the
carrier width:
\begin{equation}
e_i=\rmsn_{d_x}(x_i)P_G\in\R^{d_g},\quad P_G\in\R^{d_x\times d_g},\quad \gbar_0=0.
\end{equation}
$P_G$ has no bias. The reset gate differs from the usual sigmoid GRU reset: it uses softplus and may
therefore exceed one. With $W_r,W_z,W_h,U_r,U_z,U_h\in\R^{d_g\times d_g}$ and $b_r,b_z,b_h\in\R^{d_g}$,
\begin{align}
r_i&=\softplus(e_iW_r+\gbar_{i-1}U_r+b_r), && r_i\in(0,\infty)^{d_g},\\
z_i&=\sigma(e_iW_z+\gbar_{i-1}U_z+b_z), && z_i\in(0,1)^{d_g},\\
h_i^{\mathrm{GRU}}&=\tanh\bigl(e_iW_h+(r_i\odot\gbar_{i-1})U_h+b_h\bigr),\\
g_i^{\mathrm{raw}}&=(\one-z_i)\odot\gbar_{i-1}+z_i\odot h_i^{\mathrm{GRU}}.
\end{align}
The symbol $h_i^{\mathrm{GRU}}$ is the recurrent candidate, not the geometric-stage hidden state
introduced below.

\subsection{Signed-simplex factorization}

The raw recurrence is converted into an oriented, normalized memory representation. Let
$\beta_{\mathrm{mem}}$ and $\kappa_{\mathrm{mem}}$ denote the learnable scalars named $\beta$ and $\kappa$ in the code. Set, coordinatewise,
\begin{align}
s_i&=\tanh\bigl(\beta_{\mathrm{mem}}g_i^{\mathrm{raw}}\bigr), && s_i\in(-1,1)^{d_g},\\
m_{ij}&=\sqrt{(g_{ij}^{\mathrm{raw}})^2+\varepsilon},\\
q_{ij}&=\softplus(\kappa_{\mathrm{mem}}m_{ij})>0,\\
\alpha_{ij}&=\frac{q_{ij}}{\sum_{\ell=1}^{d_g}q_{i\ell}}, && \alpha_i\in\Delta^{d_g-1},\\
g_i&=s_i\odot\alpha_i.
\end{align}
Thus the magnitude allocation $\alpha_i$ is a simplex, whereas $g_i$ is a \emph{signed simplex-weighted
vector}; $g_i$ is not itself a probability vector. Since $|s_{ij}|<1$,
\begin{equation}
\|g_i\|_1=\sum_j|s_{ij}|\alpha_{ij}<\sum_j\alpha_{ij}=1.
\end{equation}
The entropy and corresponding effective number of active coordinates are descriptive diagnostics:
\begin{equation}
\mathcal H(\alpha_i)=-\sum_{j=1}^{d_g}\alpha_{ij}\log\alpha_{ij},\qquad
N_{\mathrm{eff}}(\alpha_i)=e^{\mathcal H(\alpha_i)}.
\end{equation}

\subsection{Carrier propagated through time}

The recurrent carrier used by the next step and by the communication attention is
\begin{equation}
\gbar_i=\gamma_c\odot\operatorname{rms}(g_i),\qquad \gamma_c\in\R^{d_g}\ \ \text{learnable and specific to the block.}
\end{equation}
The coordinatewise gain $\gamma_c$ is the learned gain of the carrier; the documented NOSCAL architecture
does not insert additional scalar gates on the two large residual paths of the block. The causal
recurrent computation at position $i$ depends on positions $\le i$ only.

% =====================================================================================================
\section{Communication attention: reading the causal carrier}\label{sec:attention}

Collect the carrier vectors into a matrix:
\begin{equation}
\Gbar=\begin{bmatrix}\gbar_1\\ \vdots\\ \gbar_L\end{bmatrix}\in\R^{L\times d_g}.
\end{equation}
Three learned projections generate multihead queries, keys and values:
\begin{equation}
Q^A=\Gbar W_Q^A,\qquad K^A=\Gbar W_K^A,\qquad V^A=\Gbar W_V^A,\quad W_Q^A,W_K^A,W_V^A\in\R^{d_g\times d_g}.
\end{equation}
Each of the $H=8$ heads has $d_h=64$ coordinates. For head $a$, use $q_i^a,k_i^a,v_i^a\in\R^{d_h}$ for
its partition of $Q^A,K^A,V^A$.

\subsection{Rotary position embedding (RoPE)}

For positions $i$, RoPE applies a block diagonal matrix containing $d_h/2$ two-dimensional rotations,
with frequencies
\begin{gather}
\omega_p=10000^{-2p/d_h},\qquad p=0,\dots,d_h/2-1,\\
R_i=\bigoplus_{p=0}^{d_h/2-1}\begin{bmatrix}\cos(i\omega_p)&-\sin(i\omega_p)\\ \sin(i\omega_p)&\cos(i\omega_p)\end{bmatrix}.
\end{gather}
Equivalently, for row-vector convention, the rotated queries and keys are written
$\tilde q_i^a=q_i^aR_i^{\T}$ and $\tilde k_i^a=k_i^aR_i^{\T}$. The rotation is applied to queries and
keys, not to values.

\subsection{Masked attention and communication residual}

For every attention head $a$, the causal scores and simplex weights are
\begin{gather}
S_{ij}^{A,a}=\begin{cases}\dfrac{\tau_{A,a}}{\sqrt{d_h}}\,\langle\tilde q_i^a,\tilde k_j^a\rangle, & j\le i,\\[4pt] -\infty, & j>i,\end{cases}\\
\rho_{ij}^{A,a}=\frac{\exp S_{ij}^{A,a}}{\sum_{\ell\le i}\exp S_{i\ell}^{A,a}},\qquad
\bigl(\rho_{i1}^{A,a},\dots,\rho_{ii}^{A,a}\bigr)\in\Delta^{i-1}.
\end{gather}
The notation $\rho^{A,a}$ is used for attention weights to avoid confusing them with the learned memory
scalar $\beta_{\mathrm{mem}}$. Here $\tau_{A,a}$ is the learned head-specific score factor.

Concatenate the eight weighted reads:
\begin{equation}
c_i^A=\operatorname*{concat}_{a=1}^{H}\Bigl(\sum_{j\le i}\rho_{ij}^{A,a}v_j^a\Bigr)\in\R^{d_g}.
\end{equation}
Communication updates the hidden representation by
\begin{equation}
x_i^{+}=x_i+c_i^AW_O,\qquad W_O\in\R^{d_g\times d_x},\qquad (W_O)_{pq}\sim\mathcal N(0,0.02^2).
\end{equation}
There is no extra learned scalar multiplier on this residual. This is the documented NOSCAL
communication path.

% =====================================================================================================
\section{Stage-wise control: a second causal attention}

The control mechanism $C$ is \emph{separate} from communication attention $A$: it has its own learned
query, key and value projections, also with $H=8$ heads of width $d_h=64$, RoPE and causal masking. It is
used to determine operator-mixture coefficients rather than the communicated residual.

Crucially, $C$ is recomputed \emph{at every geometric stage} on the currently transformed hidden
sequence. If $x_i^{(k-1)}$ is the hidden state entering stage $k$, define
\begin{equation}
h_{k,i}=\rmsn_{d_x}\Bigl(x_i^{(k-1)}\Bigr),\qquad H_k=\begin{bmatrix}h_{k,1}\\ \vdots\\ h_{k,L}\end{bmatrix}.
\end{equation}
Let the control attention have projection matrices $W_Q^C,W_K^C,W_V^C\in\R^{d_x\times d_g}$, and apply the
same masked, RoPE-based attention construction as in Section~\ref{sec:attention} to $H_k$:
\begin{equation}
Q_k^C=H_kW_Q^C,\quad K_k^C=H_kW_K^C,\quad V_k^C=H_kW_V^C,\qquad c_{k,i}^C=C(H_k)_i\in\R^{d_g}.
\end{equation}
The same \emph{control-attention module} $C$ is reused among the seven stages of a given block, although
the attention weights are recomputed at each stage. Control modules of distinct blocks are independent.
Thus the control at stage $k+1$ may read the effects of stages $1,\dots,k$.

% =====================================================================================================
\section{Lifting the current hidden state into operational space}

The initial stage state is the communicated hidden:
\begin{equation}
x_i^{(0)}=x_i^{+}.
\end{equation}
At each stage $k=1,\dots,K$, project its freshly normalized hidden state into the $D$-dimensional
operational space:
\begin{equation}
u_{k,i}=\SiLU(h_{k,i}P_g)\odot(h_{k,i}P_\up)\in\R^{D},
\end{equation}
where
\begin{equation}
P_g,P_\up\in\R^{d_x\times D},\qquad D=4096.
\end{equation}
Both projectors are \emph{shared across all seven stages of the block}. $P_g$ provides a nonlinear gate,
while $P_\up$ carries a linearly projected content vector. The Hadamard product forms the stage's
operational representation. Importantly, the lifting is performed again for every stage after the
previous update has returned to hidden space.

% =====================================================================================================
\section{Continuous local routing on a six-slot simplex}

Each stage has its own linear router head:
\begin{equation}
W_k\in\R^{d_g\times 6},\qquad b_k\in\R^{6}.
\end{equation}
From the stage-specific control read, it produces
\begin{equation}\label{eq:router}
\pi_{k,i}=\softmax\bigl(c_{k,i}^CW_k+b_k\bigr)\in\Delta^{5}.
\end{equation}
The slots are indexed as follows: $j=0$ is the \emph{null identity} (zero deformation), while
$j=1,\dots,5$ correspond to five operators of the current family's type $f_k$. For $u\in\R^D$, define
\begin{equation}
\op_{k,0}(u)=0,\qquad d_{k,i}=\sum_{j=1}^{5}\pi_{k,i,j}\,\op_{k,j}^{(f_k)}(u_{k,i})\in\R^{D}.
\end{equation}
The mixture is continuous; there is no Gumbel hard selection in this version. In the idealized boundary
case $\pi_{k,i,0}=1$, all other coefficients vanish and hence $d_{k,i}=0$ exactly. A finite softmax logit
vector typically produces strictly positive probabilities, so this identity statement is a property of
the simplex boundary, not a claim that exact one-hot output is routinely reached during training.

No single global catalogue of all operators is used. Each of the seven stage occurrences owns the
parameters of its five local operators. In particular, the three $M$ stages in~\eqref{eq:families} are
independent.

% =====================================================================================================
\section{Operator families in the lifted space}

Each operator maps $\R^D\to\R^D$ using rank-$r$ learned frames, with $D=4096$ and $r=32$. Unless
otherwise indicated, operations on intermediate $r$-vectors are elementwise. The subscripts $k,j$ on frame
matrices are temporarily suppressed for readability, but all frames belong to their \emph{specific stage
and operator slot}.

\subsection{$F$: free low-rank nonlinear field}

For the five operators of a stage of family $F$, let
\begin{equation}
A_j^F,B_j^F,U_j^F\in\R^{D\times r}.
\end{equation}
Then
\begin{equation}
\boxed{F_j(u)=\bigl[\SiLU(uA_j^F)\odot(uB_j^F)\bigr](U_j^F)^{\T}.}
\end{equation}
No fixed internal catalogue of elementary functions is imposed on the five free fields: their diversity
comes from learned frames. As $t\to0$, $\SiLU(t)=\tfrac12t+O(t^2)$, so the leading local term is
quadratic:
\begin{equation}
F_j(u)=\tfrac12\bigl[(uA_j^F)\odot(uB_j^F)\bigr](U_j^F)^{\T}+O(\|u\|^3).
\end{equation}
The documented initialization multiplies the initial $B_j^F$ frame by \texttt{F\_INIT\_GAIN}$\;=6.0$ for
empirical matching of the deformation magnitude to the other families; this scaling is an initialization
rule, not an additional runtime operation.

\subsection{$B$: tangent-to-orthogonal bending}

For each bending slot, take
\begin{equation}
T_j,N_j\in\R^{D\times r}.
\end{equation}
Let $Q_{T_j}$ be an orthonormal basis obtained from the QR factorization of $T_j$. The second frame is
orthogonalized against the span of $T_j$:
\begin{equation}
N_j^{\perp}=N_j-Q_{T_j}\bigl(Q_{T_j}^{\T}N_j\bigr).
\end{equation}
The resulting operator is
\begin{equation}
\boxed{B_j(u)=\SiLU(uT_j)(N_j^{\perp})^{\T}.}
\end{equation}
$N_j^{\perp}$ is not normalized again. If the frame is rank-deficient, the intended projection is onto the
actual column span of $T_j$; the QR expression above is the documented construction for the frame-based
implementation.

\subsection{$M$: point-dependent metric-style deformation}

For each metric-family operator, use
\begin{equation}
U_j^M,A_j^M\in\R^{D\times r},\qquad
\boxed{M_j(u)=\bigl[(uU_j^M)\odot\phi_j(uA_j^M)\bigr](U_j^M)^{\T}.}
\end{equation}
The five local nonlinearities are
\begin{equation}
\begin{aligned}
&\phi_1(q)=\tanh(q),\qquad &&\phi_2(q)=\operatorname{softsign}(q)=\frac{q}{1+|q|},\\
&\phi_3(q)=\operatorname{asinh}(q),\qquad &&\phi_4(q)=q\odot\sigma(\beta_{M,4}|q|),\\
&\phi_5(q)=\SiLU(q).
\end{aligned}
\end{equation}
The subscript on $\beta_{M,4}$ distinguishes this metric-family parameter from the memory-orientation
scalar $\beta_{\mathrm{mem}}$. The $M$ family is named for its
learned, point-dependent modifications, but the formulas alone do \emph{not} assert that a
positive-definite Riemannian metric tensor is explicitly constructed.

\subsection{$V$: oscillatory deformation family}

This family has the same low-rank frame structure as $M$:
\begin{equation}
U_j^V,A_j^V\in\R^{D\times r},\qquad
\boxed{V_j(u)=\bigl[(uU_j^V)\odot\phi_j(uA_j^V)\bigr](U_j^V)^{\T}.}
\end{equation}
The five scalar/vector-valued laws are
\begin{equation}
\begin{aligned}
&\phi_1(q)=\tanh(q), &&\phi_2(q)=\operatorname{asinh}(q),\\
&\phi_3(q)=\sin(\omega\odot q), &&\phi_4(q)=\one-\cos(\nu\odot q),\\
&\phi_5(q)=c_s\odot\sin(\omega\odot q)+c_c\odot\bigl[\one-\cos(\nu\odot q)\bigr].
\end{aligned}
\end{equation}
The necessary frequency and combination vectors $\omega,\nu,c_s,c_c\in\R^{r}$ are learned wherever the
associated law uses them. Every oscillatory operator slot has its own applicable parameters; the notation
above suppresses the slot index.

\subsection{$S$: linear shear family}

For $U_j^S,V_j^S\in\R^{D\times r}$, define
\begin{equation}
\boxed{S_j(u)=(uV_j^S)(U_j^S)^{\T}.}
\end{equation}
This is a general learned low-rank linear map of the operational vector. The term \emph{shear} is the name of this family in CASC; the learned map is not restricted by the equations to a
classical volume-preserving shear matrix.

\subsection{Family-dependent return nonlinearity and initialization}

Before being translated back to hidden space, the mixed deformation uses
\begin{equation}
\psi_F=\psi_B=\psi_M=\psi_V=\SiLU,\qquad \psi_S=\operatorname{Id}.
\end{equation}
The documented frame initialization starts from normal frame entries with variance approximately $1/D$.
Certain value-contributing frames are initialized with an additional factor $\varepsilon_0=0.02$, so typical initial relative deformation magnitudes are in the
$10^{-4}$--$10^{-3}$ range. This is an implementation initialization prescription, not a constraint on the
trained operator norms.

% =====================================================================================================
\section{Geometric cascade: composition through hidden space}

The crucial structural feature of CASC--NTP is that a deformation is \emph{immediately returned to the
hidden state} before the next deformation is evaluated. The shared return projector is
\begin{equation}
P_\dn\in\R^{D\times d_x},
\end{equation}
independent of $P_\up$; the architecture does not require $P_\dn=P_\up^{\T}$.

For every stage $k$ and token position $i$, the complete state transition is
\begin{align}
x_i^{(0)}&=x_i^{+},\\
h_{k,i}&=\rmsn_{d_x}\bigl(x_i^{(k-1)}\bigr),\\
u_{k,i}&=\SiLU(h_{k,i}P_g)\odot(h_{k,i}P_\up),\\
c_{k,i}^C&=C(H_k)_i,\\
\pi_{k,i}&=\softmax\bigl(c_{k,i}^CW_k+b_k\bigr),\\
d_{k,i}&=\sum_{j=1}^{5}\pi_{k,i,j}\,\op_{k,j}^{(f_k)}(u_{k,i}),\label{eq:mixture}\\
\Delta x_{k,i}&=\psi_{f_k}(d_{k,i})P_\dn,\\
&\boxed{x_i^{(k)}=x_i^{(k-1)}+\Delta x_{k,i}.}
\end{align}
In these equations, $H_k$ contains the normalized current states of \emph{all} sequence positions, and $C$
is causally masked. The output of the block is $B_b(x)_i=x_i^{(K)}$. This defines the deterministic
forward computation of the entire block, conditional on its parameters.

The operative cycle is
\begin{equation}
\underbrace{x^{(k-1)}}_{\text{hidden}}\;\longrightarrow\;\underbrace{h_k}_{\text{renormalize}}\;\longrightarrow\;
\underbrace{u_k}_{\text{lift}}\;\longrightarrow\;\underbrace{d_k}_{\text{routed geometry}}\;\longrightarrow\;
\underbrace{x^{(k)}}_{\text{return and add}}.
\end{equation}
The control computation is evaluated from $h_k$ at the same stage and modulates $d_k$. This is \emph{not}
a single sum of seven operators applied once to the initial $u$: both $u_k$ and $c_k^C$ change when earlier
stage updates change the hidden state.

% =====================================================================================================
\section{Full forward map and autoregressive readout}

Combining the previous sections, the causal language model is
\begin{equation}
(y_1,\dots,y_L)\xrightarrow{E}X^{[0]}\xrightarrow{B_1}\cdots\xrightarrow{B_S}X^{[S]}
\xrightarrow{\rmsn_{d_x}}H^{\mathrm{out}}\xrightarrow{E^{\T}}\Lambda\in\R^{L\times V}.
\end{equation}
For each position,
\begin{align}
h_i^{\mathrm{out}}&=\rmsn_{d_x}\bigl(x_i^{[S]}\bigr),\\
\ell_i&=h_i^{\mathrm{out}}E^{\T}\in\R^{V},\\
p_i(v)&=\frac{\exp\ell_{i,v}}{\sum_{w=0}^{V-1}\exp\ell_{i,w}},\qquad v=0,\dots,V-1.
\end{align}
The tied projection uses a \emph{single} trainable $E$ for the token embedding and the readout, rather than
maintaining independent embedding and output-head matrices.

\subsection{Causality}

Each $x_i^{[S]}$ depends only on tokens at positions at most $i$: the GRU state is causal, both attentions
use triangular masks, and the geometric operations are pointwise in position once the causal control
vectors are constructed. Therefore the logits $\ell_i$ define a next-token conditional distribution
$p_\theta(y_{i+1}\mid y_{\le i})$.

% =====================================================================================================
\section{The three probability simplices}

The architecture contains three conceptually different normalization mechanisms:
\begin{center}
\begin{tabular}{p{0.17\textwidth}p{0.17\textwidth}p{0.56\textwidth}}
\toprule
Distribution & Space & Role\\
\midrule
$\alpha_i$ & $\Delta^{d_g-1}=\Delta^{511}$ & Coordinate-weight distribution used to structure the signed recurrent memory.\\
$\rho_{i\cdot}^{A,a}$ and $\rho_{k,i\cdot}^{C,a}$ & $\Delta^{i-1}$ & Distributions over accessible past and current sequence positions for each head.\\
$\pi_{k,i}$ & $\Delta^{5}$ & Distribution over the null-identity slot and five operators of the current family.\\
\bottomrule
\end{tabular}
\end{center}
The meanings are not interchangeable: $\alpha$ distributes carrier coordinates, attention weights
distribute source positions, and $\pi$ distributes stage-level deformation choices.

% =====================================================================================================
\section{Learning objective and actual optimization mechanism}

This section distinguishes the trainable machine from the optimizer and states \emph{the training
mechanism documented for CASC--NTP}. The model is trained end-to-end using ordinary next-token
cross-entropy. For position mask $M_i\in\{0,1\}$,
\begin{equation}
\boxed{\mathcal L_{\mathrm{NTP}}(\theta)=-\frac{1}{\sum_iM_i}\sum_iM_i\log p_{\theta,i}(y_{i+1}).}
\end{equation}
Here $\theta$ comprises the embedding, recurrent memory, both attention mechanisms, local routing heads,
operator frames, projectors, and normalization parameters.

\subsection{Standard backward propagation}

Gradients are computed by \textbf{classical backpropagation}, i.e.\ reverse-mode differentiation through
the full forward computation graph (including attention, recurrent states, stage-wise geometric updates,
and tied lexical weights). With $z_{i,v}=\ell_{i,v}$ and one-hot target $e_{y_{i+1}}$,
\begin{equation}
\frac{\partial\mathcal L_{\mathrm{NTP}}}{\partial z_{i,v}}=\frac{M_i}{\sum_jM_j}\Bigl[p_{\theta,i}(v)-\one\{v=y_{i+1}\}\Bigr].
\end{equation}
For any trainable parameter $\theta_a$, the actual update gradient is
\begin{equation}
g_{t,a}=\frac{\partial\mathcal L_{\mathrm{NTP}}(\theta_t)}{\partial\theta_a},
\end{equation}
obtained via the chain rule throughout the network. In particular, the coefficients $\pi_{k,i,j}$ route
forward information and receive gradients through \eqref{eq:router}--\eqref{eq:mixture}; they \emph{do not
replace} reverse-mode derivatives as an optimization rule.

\subsection{AdamW parameter updates}

The documented trainer uses AdamW with $\beta_1=0.9$, $\beta_2=0.95$, decoupled weight decay
$\lambda=0.01$, and global gradient-norm clipping at $1.0$. For gradient $g_t$ after any prescribed
clipping, the standard AdamW equations are
\begin{align}
m_t&=\beta_1m_{t-1}+(1-\beta_1)g_t,\\
v_t&=\beta_2v_{t-1}+(1-\beta_2)(g_t\odot g_t),\\
\hat m_t&=\frac{m_t}{1-\beta_1^t},\qquad \hat v_t=\frac{v_t}{1-\beta_2^t},\\
\theta_t&=(1-\eta_t\lambda)\theta_{t-1}-\eta_t\frac{\hat m_t}{\sqrt{\hat v_t}+\varepsilon_{\mathrm{AdamW}}}.
\end{align}
These are conventional optimizer equations supplied here to identify the training mechanism;
$\varepsilon_{\mathrm{AdamW}}$ is the optimizer stabilization constant and is \emph{not} the
$\varepsilon=10^{-6}$ used in the geometric normalizations. The released trainer uses the PyTorch default $\varepsilon_{\mathrm{AdamW}}=10^{-8}$. The
learning-rate schedule belongs to the training protocol, not to the mathematical definition of the machine.

The published run, NTP\_A\_4B, used an initial learning rate of $6\times10^{-4}$ after a warmup of 2000
steps and a warmup--stable--decay schedule with a linear decay over the last 15\% of training. It completed
49{,}215 optimizer steps over $3.592\times10^{9}$ prediction tokens, one full pass over its corpus. These
numbers document that run; they are not requirements for constructing or evaluating CASC--NTP.

% =====================================================================================================
\section{Parameter count and sharing rules}

The full-size configuration has 175{,}218{,}812 trainable parameters, partitioned as follows:
\begin{center}
\begin{tabular}{lr}
\toprule
Component & Trainable parameters\\
\midrule
Tied embedding $E$ & 65.536\,M\\
Causal memory and attention mechanisms & 19.931\,M\\
Shared within-block operational projectors & 50.332\,M\\
Geometric operator-family frames & 39.325\,M\\
Stage routers and normalization scales & 0.095\,M\\
\midrule
\textbf{Total} & $\mathbf{175.219}$\,\textbf{M}\\
\bottomrule
\end{tabular}
\end{center}
The three large projection matrices $P_g,P_\up,P_\dn$ contribute, per block,
\begin{equation}
3d_xD=3(1024)(4096)=12{,}582{,}912
\end{equation}
weights. Across four blocks, this equals 50{,}331{,}648. The tied embedding has $Vd_x=65{,}536{,}000$
parameters. Component totals are rounded to three decimal places in millions.

Parameter sharing is precisely localized: the attention $C$ and operational projectors are reused across
stages of a \emph{single} block; operator-frame parameters and router head parameters are distinct for each
stage; blocks are separate; and lexical weights are tied between input embedding and final readout.

% =====================================================================================================
\section{Architectural diagnostics and interpretation}

The trainer logs the following diagnostics per block and geometric stage, without
treating them as extra losses:
\begin{equation}
\pi_{k,i,0},\qquad \exp\mathcal H(\pi_{k,i}),\qquad \frac{\|d_{k,i}\|_2}{\|u_{k,i}\|_2},\qquad
\frac{\|\Delta x_{k,i}\|_2}{\|x_i^{(k-1)}\|_2}.
\end{equation}
For the cumulative dynamics, relevant quantities include
\begin{equation}
\cos(\Delta x_{k,i},\Delta x_{\ell,i}),\qquad \frac{\bigl\|\sum_{j\le k}\Delta x_{j,i}\bigr\|_2}{\|x_i^{+}\|_2},\qquad
\frac{\|u_{k,i}-u_{k-1,i}\|_2}{\|u_{k-1,i}\|_2},\qquad \cos(u_{k,i},u_{k-1,i}),
\end{equation}
where cosine similarity is defined when its arguments have nonzero norm, and denominators should be
numerically safeguarded when these diagnostics are implemented. Percentile-based statistics of
deformation magnitudes (e.g.\ $r_{90}$ and $r_{99}$ in the training log) and norms of the $F$ family's
intermediate factors are also useful to identify concentration, cancellation, stage inactivity, or
uncontrolled initialization growth.

\subsection{Meaning and limits of the geometric vocabulary}

The terms \emph{geometric deformation}, \emph{metric}, \emph{bending}, \emph{shear}, and \emph{manifold}
describe the architecture's intended interpretation: learned operators move a hidden representation
through a sequence of locally defined transformations, and each update changes the inputs to the
subsequent stage. The equations establish this \emph{computational} interpretation. They do not, by
themselves, prove that the hidden-state data occupy a smooth embedded submanifold, that any learned update
is a diffeomorphism, that all mappings are invertible, or that the $M$ family defines a positive-definite
metric. Those would require additional mathematical assumptions and separate verification. In particular,
the low-rank, gated, residual maps are not constrained here to preserve dimension, distance, topology,
volume, or curvature in a specified formal sense.

% =====================================================================================================
\section{Compact specification of the complete machine}

For $b=1,\dots,S$, the block acts as
\begin{equation}
\boxed{
\begin{array}{c}
\text{recurrent signed-simplex carrier }\Gbar^{(b)}\xrightarrow{A^{(b)}}c^{A,(b)}\;\longrightarrow\;x^{+,(b)},\\[3pt]
x^{(0,b)}=x^{+,(b)},\\[3pt]
\left.\begin{aligned}
h_k^{(b)}&=\rmsn\bigl(x^{(k-1,b)}\bigr),\\
u_k^{(b)}&=\SiLU\bigl(h_k^{(b)}P_g^{(b)}\bigr)\odot\bigl(h_k^{(b)}P_\up^{(b)}\bigr),\\
c_k^{C,(b)}&=C^{(b)}\bigl(h_k^{(b)}\bigr),\\
\pi_k^{(b)}&=\softmax\bigl(c_k^{C,(b)}W_k^{(b)}+b_k^{(b)}\bigr),\\
d_k^{(b)}&=\textstyle\sum_{j=1}^{5}\pi_{k,j}^{(b)}\,\op_{k,j}^{(f_k,b)}\bigl(u_k^{(b)}\bigr),\\
x^{(k,b)}&=x^{(k-1,b)}+\psi_{f_k}\bigl(d_k^{(b)}\bigr)P_\dn^{(b)}
\end{aligned}\right\}\quad k=1,\dots,7.
\end{array}}
\end{equation}
All operations in this display are performed over token positions as described earlier, with causal
dependencies preserved by the recurrent update and masked attentions. Finally,
\begin{equation}
\boxed{
\begin{aligned}
X^{[0]}&=E[y_1,\dots,y_L],\\
X^{[S]}&=B_S\circ\cdots\circ B_1\bigl(X^{[0]}\bigr),\\
\Lambda&=\rmsn_{d_x}\bigl(X^{[S]}\bigr)E^{\T},\\
p_\theta(y_{i+1}\mid y_{\le i})&=\softmax(\Lambda_i)_{y_{i+1}},\\
g_\theta&=\nabla_\theta\mathcal L_{\mathrm{NTP}}\quad\text{(classical backpropagation)},\\
\theta_{t+1}&=\operatorname{AdamW}(\theta_t,g_\theta).
\end{aligned}}
\end{equation}

% =====================================================================================================
\section*{Verification against the released code}
\addcontentsline{toc}{section}{Verification against the released code}

The statements of this document were checked against the released source code (\texttt{casc.py},
\texttt{train.py}) with \texttt{test\_casc.py}:
\begin{itemize}
\item tensor shapes and the parameter count: $175{,}218{,}812$ trainable parameters in 274 state-dict
entries, whose names and shapes are exactly those of the published checkpoint (loaded with
\texttt{strict=True});
\item parameter sharing: one control attention $C$ and one projector triplet $(P_g,P_\up,P_\dn)$ per
block, independent operator frames and router heads for every stage, a single tied matrix $E$;
\item the equations of Sections~3--10, by independent re-implementation: signed-simplex GRU, causal
attentions, identity slot, the $F$, $B$, $M$, $V$, $S$ operators, $\psi$ per family, renormalization and
control recomputed at every stage, causality of the whole machine;
\item the optimizer: \texttt{torch.optim.AdamW} with $\beta_1=0.9$, $\beta_2=0.95$, $\lambda=0.01$ and the
PyTorch default $\varepsilon_{\mathrm{AdamW}}=10^{-8}$, global gradient-norm clipping at $1.0$.
\end{itemize}
The measurements of the published run are reported in \texttt{report.html}.

\end{document}
